
\documentclass[11pt,a4paper]{article}
\usepackage{amsmath,amssymb}
\usepackage{fontspec}
\usepackage{xeCJK}
\usepackage{geometry}
\usepackage{xcolor}
\usepackage{url}
\usepackage{hyperref}
\usepackage{fvextra}
\usepackage{titlesec}
\usepackage{fancyhdr}
\usepackage{enumitem}

% Long URLs and identifiers must be able to break across lines.
\tolerance=2000
\emergencystretch=3em
\hbadness=10000

\defaultfontfeatures{Ligatures=TeX}
\setmainfont{Times New Roman}
\setsansfont{Helvetica Neue}
\setmonofont{Menlo}[Scale=0.88]
\setCJKmainfont{PingFang SC}
\setCJKsansfont{PingFang SC}
\setCJKmonofont{PingFang SC}
\xeCJKsetup{CJKmath=true}

\definecolor{zhcolor}{RGB}{26,26,26}
\definecolor{linkcolor}{RGB}{26,92,158}
\definecolor{codebg}{RGB}{244,245,247}
\definecolor{codeframe}{RGB}{216,220,226}
\hypersetup{colorlinks=true,linkcolor=linkcolor,urlcolor=linkcolor,
            citecolor=linkcolor,breaklinks=true}
\renewcommand{\UrlBreaks}{\do\/\do\-\do\.\do\_\do\?\do\&\do\+\do\=\do\#}

\geometry{top=2.2cm,bottom=2.2cm,left=2.3cm,right=2.3cm}
\sloppy

\titleformat{\section}{\Large\bfseries\sffamily}{}{0em}{}
\titleformat{\subsection}{\large\bfseries\sffamily}{}{0em}{}

\pagestyle{fancy}
\fancyhf{}
\fancyhead[L]{\small\sffamily tiny-llm · Week 1 Day 3--5 双语对照}
\fancyhead[R]{\small\sffamily\thepage}
\renewcommand{\headrulewidth}{0.4pt}

\newenvironment{zhpar}{%
  \par\smallskip\begingroup\color{zhcolor}\small\ignorespaces%
}{\endgroup\par\smallskip}

\DefineVerbatimEnvironment{codeblock}{Verbatim}{
  frame=single,framesep=5pt,rulecolor=\color{codeframe},
  fillcolor=\color{codebg},fontsize=\small,formatcom=\color{black},
  baselinestretch=1.0,xleftmargin=2pt,breaklines=true,
  breakanywhere=true
}

\begin{document}
\begin{center}
{\LARGE\bfseries\sffamily tiny-llm: Week 1 Day 3--5}\\[0.4em]
{\large\sffamily Grouped Query Attention · RMSNorm 与 MLP · Qwen3 模型}\\[0.3em]
{\small\sffamily 中英双语对照 · Bilingual Edition}
\end{center}
\vspace{0.6em}
\section*{Day 3: Grouped Query Attention (GQA)}

On Day 3, we will implement grouped-query attention (GQA). Qwen3 uses GQA to reduce the computational and memory costs of the key (K) and value (V) projections. In multi-head attention (MHA), every query (Q) head has a corresponding K and V head. With GQA, groups of Q heads share K and V heads. Multi-query attention (MQA) is the special case in which every Q head shares a single K/V head pair.

\begin{zhpar}
第 3 天，我们将实现分组查询注意力（Grouped-Query Attention，GQA）。Qwen3 使用 GQA 来降低键（K）与值（V）投影的计算与内存开销。在多头注意力（MHA）中，每个查询（Q）头都有与之对应的 K 头和 V 头；而在 GQA 中，若干个 Q 头组成一组，共享同一组 K/V 头。多查询注意力（MQA）是 GQA 的特例：所有 Q 头共享唯一的一对 K/V 头。
\end{zhpar}
\textbf{延伸阅读 / Readings}

\begin{itemize}[leftmargin=1.6em,itemsep=0.15em,parsep=0pt]
\item \href{https://arxiv.org/abs/2305.13245}{GQA paper}
\item \href{https://github.com/ml-explore/mlx-lm/blob/main/mlx_lm/models/qwen3.py}{Qwen3 layers in mlx-lm}
\item \href{https://pytorch.org/docs/stable/generated/torch.nn.functional.scaled_dot_product_attention.html}{PyTorch scaled dot-product attention with \texttt{enable\_gqa=True}}
\item \href{https://pytorch.org/torchtune/0.3/generated/torchtune.modules.MultiHeadAttention.html}{\texttt{torchtune.modules.MultiHeadAttention}}
\end{itemize}

\subsection*{Task 1: Implement \texttt{scaled\_dot\_product\_attention\_grouped}}

\subsection*{任务 1：实现 \texttt{scaled\_dot\_product\_attention\_grouped}}

You will need to modify the following file:

\begin{zhpar}
你需要修改以下文件：
\end{zhpar}
\begin{codeblock}
src/tiny_llm/attention.py
\end{codeblock}

In this task, we will implement grouped scaled dot-product attention, which forms the core of GQA.

\begin{zhpar}
本任务将实现“分组缩放点积注意力”，它是 GQA 的核心。
\end{zhpar}
Implement \texttt{scaled\_dot\_product\_attention\_grouped} in \texttt{src/tiny\_llm/attention.py}. It is similar to standard scaled dot-product attention, but it supports a number of query heads that is a multiple of the number of key/value heads.

\begin{zhpar}
请在 \texttt{src/tiny\_llm/attention.py} 中实现 \texttt{scaled\_dot\_product\_attention\_grouped}。它与标准的缩放点积注意力类似，区别在于：查询头的数量可以是键/值头数量的整数倍。
\end{zhpar}
The main process is the same as standard scaled dot-product attention. The difference is that K and V heads are shared across multiple Q heads. Instead of \texttt{H\_q} separate K and V heads, there are \texttt{H} K and V heads, each shared by \texttt{n\_repeats = H\_q // H} query heads.

\begin{zhpar}
整体流程与标准缩放点积注意力一致，区别在于 K/V 头被多个 Q 头共享：不存在 \texttt{H\_q} 个彼此独立的 K/V 头，而是只有 \texttt{H} 个 K/V 头，每个被 \texttt{n\_repeats = H\_q // H} 个查询头共享。
\end{zhpar}
Reshape \texttt{query}, \texttt{key}, and \texttt{value} so that K and V can be broadcast to the query heads in their respective groups during the matrix multiplications.

\begin{zhpar}
请调整 \texttt{query}、\texttt{key}、\texttt{value} 的形状，使 K 和 V 能在矩阵乘法中广播到各自分组内的所有查询头上。
\end{zhpar}
\begin{itemize}[leftmargin=1.6em,itemsep=0.15em,parsep=0pt]
\item Separate the \texttt{H} and \texttt{n\_repeats} dimensions in \texttt{query}.
\item Add a dimension of size 1 for \texttt{n\_repeats} in \texttt{key} and \texttt{value} so that they broadcast across each group.
\end{itemize}

\begin{zhpar}
\begin{itemize}[leftmargin=1.6em,itemsep=0.15em,parsep=0pt]
\item 在 \texttt{query} 中把 \texttt{H} 与 \texttt{n\_repeats} 两个维度拆开；
\item 在 \texttt{key} 和 \texttt{value} 中插入一个大小为 1 的 \texttt{n\_repeats} 维度，使其能在每个分组内广播。
\end{itemize}
\end{zhpar}
Then perform scaled dot-product attention: matrix multiplication, scaling, optional masking, softmax, and a final matrix multiplication. Broadcasting handles the head sharing without materializing repeated K and V tensors.

\begin{zhpar}
随后执行缩放点积注意力的标准步骤：矩阵乘法、缩放、可选的掩码、softmax，以及最后一次矩阵乘法。借助广播即可实现头共享，无需真正把 K、V 复制扩展成重复的大张量。
\end{zhpar}
Using broadcasting instead of repeating K and V is more efficient because it avoids creating copies of the same data.

\begin{zhpar}
相比“重复复制”K 和 V，使用广播更高效，因为它避免了创建同一份数据的多个副本。
\end{zhpar}
Finally, reshape the result to the expected output shape.

\begin{zhpar}
最后，把结果 reshape 成期望的输出形状。
\end{zhpar}
\begin{codeblock}
N.. is zero or more dimensions for batches
H_q is the number of query heads
H is the number of key/value heads (H_q must be divisible by H)
L is the query sequence length
S is the key/value sequence length
D is the head dimension

query: N.. x H_q x L x D
key: N.. x H x S x D
value: N.. x H x S x D
mask: N.. x H_q x L x S
output: N.. x H_q x L x D
\end{codeblock}

In addition to grouped heads, this function supports different query and key/value sequence lengths: Q uses length \texttt{L}, while K and V use length \texttt{S}.

\begin{zhpar}
除了分组头之外，该函数还支持查询与键/值使用不同的序列长度：Q 的长度为 \texttt{L}，而 K、V 的长度为 \texttt{S}。
\end{zhpar}
You can test your implementation by running the following command:

\begin{zhpar}
可以运行以下命令测试你的实现：
\end{zhpar}
\begin{codeblock}
pdm run test --week 1 --day 3 -- -k task_1
\end{codeblock}

\subsection*{Task 2: Causal Masking}

\subsection*{任务 2：因果掩码}

\textbf{延伸阅读 / Readings}

\begin{itemize}[leftmargin=1.6em,itemsep=0.15em,parsep=0pt]
\item \href{https://www.gilesthomas.com/2025/03/llm-from-scratch-9-causal-attention}{Writing an LLM from scratch, part 9 -- causal attention}
\end{itemize}

In this task, we will add causal masking to grouped attention.

\begin{zhpar}
本任务将为分组注意力加上因果掩码（causal masking）。
\end{zhpar}
Causal masking prevents attention from reading future tokens. When \texttt{mask} is set to the string \texttt{"causal"}, apply a causal mask.

\begin{zhpar}
因果掩码用于阻止注意力读取未来的 token。当 \texttt{mask} 被设为字符串 \texttt{"causal"} 时，就应用因果掩码。
\end{zhpar}
The additive causal mask has shape \texttt{(L, S)}, where \texttt{L} is the query sequence length and \texttt{S} is the key/value sequence length. Allowed positions contain 0, and masked positions contain \texttt{-inf}. When \texttt{S} is greater than \texttt{L}, shift the diagonal by \texttt{S - L} so that the queries correspond to the final \texttt{L} positions in the key/value sequence. For example, if \texttt{L = 3} and \texttt{S = 5}, the mask is:

\begin{zhpar}
这个加性（additive）因果掩码的形状为 \texttt{(L, S)}，其中 \texttt{L} 是查询序列长度，\texttt{S} 是键/值序列长度。允许 attend 的位置取值为 0，被屏蔽的位置取值为 \texttt{-inf}。当 \texttt{S} 大于 \texttt{L} 时，需要把对角线偏移 \texttt{S - L}，使查询对应键/值序列中最后的 \texttt{L} 个位置。例如当 \texttt{L = 3}、\texttt{S = 5} 时，掩码为：
\end{zhpar}
\begin{codeblock}
0   0   0   -inf -inf
0   0   0   0    -inf
0   0   0   0    0
\end{codeblock}

Implement \texttt{causal\_mask} in \texttt{src/tiny\_llm/attention.py}, then use it in \texttt{scaled\_dot\_product\_attention\_grouped}. Note that our shifted diagonal for \texttt{L != S} differs from the default behavior of some attention APIs.

\begin{zhpar}
请在 \texttt{src/tiny\_llm/attention.py} 中实现 \texttt{causal\_mask}，并在 \texttt{scaled\_dot\_product\_attention\_grouped} 中使用它。注意：我们在 \texttt{L != S} 时采用的对角线偏移方式，与某些注意力 API 的默认行为不同。
\end{zhpar}
You can test your implementation by running the following command:

\begin{zhpar}
可以运行以下命令测试你的实现：
\end{zhpar}
\begin{codeblock}
pdm run test --week 1 --day 3 -- -k task_2
\end{codeblock}

\subsection*{Task 3: Qwen3 Grouped Query Attention}

\subsection*{任务 3：Qwen3 的分组查询注意力}

In this task, we will implement Qwen3's grouped-query attention. Modify the following file:

\begin{zhpar}
本任务将实现 Qwen3 的分组查询注意力。请修改以下文件：
\end{zhpar}
\begin{codeblock}
src/tiny_llm/qwen3_week1.py
\end{codeblock}

\texttt{Qwen3MultiHeadAttention} implements attention for Qwen3. Follow this pseudocode:

\begin{zhpar}
\texttt{Qwen3MultiHeadAttention} 实现了 Qwen3 的注意力。请按如下伪代码实现：
\end{zhpar}
\begin{codeblock}
x: B, L, E
q = linear(x, wq) -> B, L, H_q, D
k = linear(x, wk) -> B, L, H, D
v = linear(x, wv) -> B, L, H, D
q = rms_norm(q, q_norm)
k = rms_norm(k, k_norm)
q = rope(q, offset=slice(0, L))
k = rope(k, offset=slice(0, L))
(transpose as needed)
x = scaled_dot_product_attention_grouped(q, k, v, scale, mask) -> B, H_q, L, D  # use float32
(transpose as needed)
x = linear(x, wo) -> B, L, E
\end{codeblock}

Qwen3 attention has no Q/K/V projection biases, and it applies RMSNorm to each Q and K head before RoPE. We will implement the reusable \texttt{RMSNorm} layer on Day 4, so call \texttt{mx.fast.rms\_norm} directly for \texttt{q\_norm} and \texttt{k\_norm} today. Use non-traditional RoPE.

\begin{zhpar}
Qwen3 的注意力中 Q/K/V 投影都不带偏置，并且在 RoPE 之前对每个 Q 头和 K 头分别做 RMSNorm。可复用的 \texttt{RMSNorm} 层会在第 4 天实现，因此今天请直接调用 \texttt{mx.fast.rms\_norm} 来完成 \texttt{q\_norm} 和 \texttt{k\_norm}。RoPE 使用非传统（non-traditional）布局。
\end{zhpar}
You can test your implementation by running the following command:

\begin{zhpar}
可以运行以下命令测试你的实现：
\end{zhpar}
\begin{codeblock}
pdm run test --week 1 --day 3 -- -k task_3
\end{codeblock}

At the end of the day, you should be able to pass all tests of this day:

\begin{zhpar}
在今天结束时，你应该能通过当天的全部测试：
\end{zhpar}
\begin{codeblock}
pdm run test --week 1 --day 3
\end{codeblock}

\clearpage
\section*{Day 4: RMSNorm and the Multilayer Perceptron}

\section*{第 4 天：RMSNorm 与多层感知机}

On Day 4, we will implement two important components of the Qwen3 Transformer architecture: RMSNorm and the multilayer perceptron (MLP), also known as the feed-forward network. RMSNorm is a normalization technique with less computational overhead than traditional layer normalization. The MLP applies nonlinear transformations after the attention block.

\begin{zhpar}
第 4 天，我们将实现 Qwen3 Transformer 架构中的两个重要组件：RMSNorm 与多层感知机（MLP，也称前馈网络）。RMSNorm 是一种归一化方法，计算开销比传统的层归一化（LayerNorm）更低。MLP 则在注意力块之后施加非线性变换。
\end{zhpar}
\subsection*{Task 1: Implement \texttt{RMSNorm}}

\subsection*{任务 1：实现 \texttt{RMSNorm}}

In this task, we will implement the \texttt{RMSNorm} layer.

\begin{zhpar}
本任务将实现 \texttt{RMSNorm} 层。
\end{zhpar}
\begin{codeblock}
src/tiny_llm/layer_norm.py
\end{codeblock}

Day 3 used \texttt{mx.fast.rms\_norm} directly so that the GQA chapter could stay focused on attention. This task implements the same normalization rule as a reusable layer. From this point on, the Transformer block, final model normalization, and Q/K normalization path can use your \texttt{RMSNorm} implementation.

\begin{zhpar}
第 3 天直接调用了 \texttt{mx.fast.rms\_norm}，是为了让 GQA 一章专注在注意力本身。本任务则把同一套归一化规则实现成一个可复用的层。从现在起，Transformer 块、模型最后的归一化，以及 Q/K 归一化路径，都可以改用你自己的 \texttt{RMSNorm} 实现。
\end{zhpar}
\textbf{ 延伸阅读 / Readings}

\begin{itemize}[leftmargin=1.6em,itemsep=0.15em,parsep=0pt]
\item \href{https://arxiv.org/abs/1910.07467}{Root Mean Square Layer Normalization}
\item \href{https://github.com/ml-explore/mlx-lm/blob/main/mlx_lm/models/qwen3.py}{Qwen3 layers implementation in mlx-lm (includes RMSNorm)} - See \texttt{RMSNorm}.
\end{itemize}

RMSNorm is defined as:

\begin{zhpar}
RMSNorm 的定义如下：
\end{zhpar}
\[
y = \frac{x}{\sqrt{\text{mean}(x^2) + \epsilon}} \cdot \text{weight}
\]

where:

\begin{zhpar}
其中：
\end{zhpar}
\begin{itemize}[leftmargin=1.6em,itemsep=0.15em,parsep=0pt]
\item \texttt{x} is the input tensor.
\item \texttt{weight} is a learned scaling parameter.
\item \texttt{epsilon} (\texttt{eps}) is a small constant, such as \texttt{1e-5} or \texttt{1e-6}, added for numerical stability.
\item \texttt{mean(x\textasciicircum{}2)} is the mean of the squared elements along the final dimension.
\end{itemize}

\begin{zhpar}
\begin{itemize}[leftmargin=1.6em,itemsep=0.15em,parsep=0pt]
\item \texttt{x} 是输入张量；
\item \texttt{weight} 是可学习的缩放参数；
\item \texttt{epsilon}（\texttt{eps}）是为了数值稳定而加入的极小常数，例如 \texttt{1e-5} 或 \texttt{1e-6}；
\item \texttt{mean(x\textasciicircum{}2)} 是沿最后一个维度求得的平方元素均值。
\end{itemize}
\end{zhpar}
Apply normalization independently to each feature vector along the input's final dimension. Cast the input to \texttt{float32} for the normalization calculation, including the mean, to preserve precision when the original values use \texttt{float16} or \texttt{bfloat16}. Cast the normalized value back to the input dtype before applying \texttt{weight}. This matches the low-precision path used by MLX's fast RMSNorm kernels: normalization statistics are accumulated in \texttt{float32}, while the final scaling happens in the model dtype.

\begin{zhpar}
请沿着输入的最后一个维度，对每个特征向量独立地做归一化。做归一化计算（包括求均值）时，先把输入转成 \texttt{float32}，以在原始值为 \texttt{float16} 或 \texttt{bfloat16} 时保住精度；在乘上 \texttt{weight} 之前，再把归一化后的值转回输入原始 dtype。这与 MLX 快速 RMSNorm 内核的低精度路径一致：归一化统计量在 \texttt{float32} 中累加，而最后的缩放仍在模型 dtype 下进行。
\end{zhpar}
\begin{codeblock}
D is the embedding dimension.

x: N.. x D
weight: D
output: N.. x D
\end{codeblock}

You can test your implementation by running:

\begin{zhpar}
可以运行以下命令测试你的实现：
\end{zhpar}
\begin{codeblock}
pdm run test --week 1 --day 4 -- -k task_1
\end{codeblock}

\subsection*{Task 2: Implement the MLP Block}

\subsection*{任务 2：实现 MLP 块}

In this task, we will implement the MLP block named \texttt{Qwen3MLP}.

\begin{zhpar}
本任务将实现名为 \texttt{Qwen3MLP} 的 MLP 块。
\end{zhpar}
\begin{codeblock}
src/tiny_llm/qwen3_week1.py
\end{codeblock}

The original Transformer uses a simple position-wise feed-forward network (FFN) in each block. It consists of two linear transformations with a ReLU activation between them.

\begin{zhpar}
最初的 Transformer 在每个块中使用简单的逐位置前馈网络（FFN），它由两个线性变换及其之间的 ReLU 激活组成。
\end{zhpar}
Modern Transformer architectures, including Qwen3, often use more advanced FFN variants. Qwen3 uses SwiGLU, a gated linear unit (GLU) variant.

\begin{zhpar}
包括 Qwen3 在内的现代 Transformer 架构通常使用更先进的 FFN 变体。Qwen3 使用 SwiGLU，一种门控线性单元（GLU）变体。
\end{zhpar}
A plain FFN can be abstracted as:

\begin{zhpar}
一个普通 FFN 可以抽象为：
\end{zhpar}
\begin{codeblock}
h = activation(W_up(x))
out = W_down(h)
\end{codeblock}

A GLU keeps the same expand-then-project-back shape but adds another projection that gates the intermediate features before \texttt{W\_down}. This gives the MLP a learned, input-dependent way to control which intermediate channels matter, rather than applying an activation only to the features produced by \texttt{W\_up}.

\begin{zhpar}
GLU 保留了同样的“先升维、再投影回去”的结构，但在 \texttt{W\_down} 之前额外加入一个投影，用来对中间特征做门控。这让 MLP 拥有一种可学习、依赖输入的方式来控制哪些中间通道更重要，而不仅仅是对 \texttt{W\_up} 产生的特征施加一个激活函数。
\end{zhpar}
SwiGLU is the GLU variant used by Qwen3:

\begin{zhpar}
SwiGLU 即 Qwen3 采用的 GLU 变体：
\end{zhpar}
\begin{codeblock}
u = W_up(x)
g = SiLU(W_gate(x))
out = W_down(g * u)
\end{codeblock}

\textbf{ 延伸阅读 / Readings}

\begin{itemize}[leftmargin=1.6em,itemsep=0.15em,parsep=0pt]
\item \href{https://arxiv.org/abs/1706.03762}{Attention is All You Need (Transformer Paper, Section 3.3 "Position-wise Feed-Forward Networks")}
\item \href{https://arxiv.org/pdf/1612.08083}{GLU paper: Language Modeling with Gated Convolutional Networks}
\item \href{https://arxiv.org/pdf/1710.05941}{SiLU (Swish) activation function}
\item \href{https://arxiv.org/abs/2002.05202v1}{SwiGLU paper: GLU Variants Improve Transformer}
\item \href{https://pytorch.org/docs/stable/generated/torch.nn.SiLU.html}{PyTorch SiLU documentation}
\item \href{https://github.com/ml-explore/mlx-lm/blob/main/mlx_lm/models/qwen3.py}{Qwen3 layers implementation in mlx-lm (includes MLP)}
\end{itemize}

SwiGLU combines a GLU with the SiLU (sigmoid linear unit) activation function:

\begin{zhpar}
SwiGLU 把 GLU 与 SiLU（sigmoid linear unit）激活函数结合起来：
\end{zhpar}
\begin{itemize}[leftmargin=1.6em,itemsep=0.15em,parsep=0pt]
\item A GLU gates one linear projection of the input with another, using element-wise multiplication to control which features pass through.
\item SiLU is a smooth, non-monotonic activation function. Unlike ReLU, it has no zero-gradient region across all negative inputs and can produce nonzero outputs for negative values.
\end{itemize}

\begin{zhpar}
\begin{itemize}[leftmargin=1.6em,itemsep=0.15em,parsep=0pt]
\item GLU 用输入的另一个线性投影去门控其中一个线性投影，通过逐元素乘法控制哪些特征得以通过；
\item SiLU 是平滑且非单调的激活函数。与 ReLU 不同，它在所有负输入上并不存在梯度恒为零的区域，并且对负值也能产生非零输出。
\end{itemize}
\end{zhpar}
First, implement \texttt{silu} in \texttt{basics.py}. It takes a tensor of shape \texttt{N.. x I} and returns a tensor with the same shape:

\begin{zhpar}
首先，在 \texttt{basics.py} 中实现 \texttt{silu}。它接收形状为 \texttt{N.. x I} 的张量，并返回同形状张量：
\end{zhpar}
\[
\text{SiLU}(x) = x * \text{sigmoid}(x) = \frac{x}{1 + e^{-x}}
\]

Compute the sigmoid part in a numerically stable way:

\begin{zhpar}
请以数值稳定的方式计算 sigmoid 部分：
\end{zhpar}
\begin{codeblock}
if x >= 0:
    sigmoid(x) = 1 / (1 + exp(-x))
else:
    sigmoid(x) = exp(x) / (1 + exp(x))
\end{codeblock}

The negative branch is algebraically equivalent to the direct sigmoid formula, but it prevents \texttt{exp(-x)} from becoming \texttt{exp(large positive)} when \texttt{x} is a large negative value. In vector code, first compute \texttt{z = exp(-abs(x))}. Use \texttt{z / (1 + z)} for negative inputs and \texttt{1 / (1 + z)} otherwise. Do not rewrite the negative branch as \texttt{1 - 1 / (1 + z)}: in low precision, the fraction can round to \texttt{1}, and the subtraction then incorrectly produces zero.

\begin{zhpar}
负半轴分支在代数上与 sigmoid 的直接公式等价，但它避免了当 \texttt{x} 为很大负数时 \texttt{exp(-x)} 变成 \texttt{exp(很大的正数)} 而溢出。在向量化代码中，先计算 \texttt{z = exp(-abs(x))}：负输入用 \texttt{z / (1 + z)}，其余用 \texttt{1 / (1 + z)}。不要把负半轴分支改写成 \texttt{1 - 1 / (1 + z)}：在低精度下这个分式可能舍入到 \texttt{1}，随后的减法就会错误地得到 0。
\end{zhpar}
Then implement \texttt{Qwen3MLP}. Qwen3's MLP contains:

\begin{zhpar}
接着实现 \texttt{Qwen3MLP}。Qwen3 的 MLP 包含：
\end{zhpar}
\begin{itemize}[leftmargin=1.6em,itemsep=0.15em,parsep=0pt]
\item A gate projection ($W_{gate}$)
\item An up projection ($W_{up}$)
\item SiLU applied to the gate projection's output
\item An element-wise product of the activated gate output and the up-projection output
\item A final down projection ($W_{down}$)
\end{itemize}

\begin{zhpar}
\begin{itemize}[leftmargin=1.6em,itemsep=0.15em,parsep=0pt]
\item 门控投影（$W_{gate}$）
\item 升维投影（$W_{up}$）
\item 对门控投影输出施加 SiLU
\item 将激活后的门控输出与升维投影输出逐元素相乘
\item 最后的降维投影（$W_{down}$）
\end{itemize}
\end{zhpar}
This can be expressed as:

\begin{zhpar}
可以表示为：
\end{zhpar}
\[
\text{MLP}(x) = W_{down}(\text{SiLU}(W_{gate}(x)) \odot W_{up}(x))
\]

where $\odot$ denotes element-wise multiplication. Qwen3's MLP projections do not use biases.

\begin{zhpar}
其中 $\odot$ 表示逐元素相乘。Qwen3 的 MLP 投影均不使用偏置。
\end{zhpar}
\begin{codeblock}
N.. is zero or more dimensions for batches
E is hidden_size (embedding dimension of the model)
I is intermediate_size (dimension of the hidden layer in MLP)
L is the sequence length

input: N.. x L x E
w_gate: I x E
w_up: I x E
w_down: E x I
output: N.. x L x E
\end{codeblock}

You can test your implementation by running:

\begin{zhpar}
可以运行以下命令测试你的实现：
\end{zhpar}
\begin{codeblock}
pdm run test --week 1 --day 4 -- -k task_2
\end{codeblock}

At the end of the day, you should be able to pass all tests of this day:

\begin{zhpar}
在今天结束时，你应该能通过当天的全部测试：
\end{zhpar}
\begin{codeblock}
pdm run test --week 1 --day 4
\end{codeblock}

\clearpage
\section*{Day 5: The Qwen3 Model}

\section*{第 5 天：Qwen3 模型}

On Day 5, we will combine the components from the previous chapters into the complete Qwen3 model.

\begin{zhpar}
第 5 天，我们将把前几章实现的各个组件组合成完整的 Qwen3 模型。
\end{zhpar}
Model-level tests require the corresponding model files. Start with the default 0.6B model; download the larger models only if you want to test them as well:

\begin{zhpar}
模型级别的测试需要对应的模型文件。请先从默认的 0.6B 模型开始；更大的模型只在你也想测试它们时才下载：
\end{zhpar}
\begin{codeblock}
hf download Qwen/Qwen3-0.6B-MLX-4bit
# Optional larger models:
hf download Qwen/Qwen3-1.7B-MLX-4bit
hf download Qwen/Qwen3-4B-MLX-4bit
\end{codeblock}

Tests that require an unavailable model will be skipped.

\begin{zhpar}
需要某个未下载模型的测试会被自动跳过。
\end{zhpar}
\subsection*{Task 1: Implement \texttt{Qwen3TransformerBlock}}

\subsection*{任务 1：实现 \texttt{Qwen3TransformerBlock}}

\begin{codeblock}
src/tiny_llm/qwen3_week1.py
\end{codeblock}

\textbf{ 延伸阅读 / Readings}

\begin{itemize}[leftmargin=1.6em,itemsep=0.15em,parsep=0pt]
\item \href{https://medium.com/@akhileshkapse/a-simplified-explanation-of-the-transformer-block-must-read-blog-for-nlp-enthusiasts-12ef240a62ac}{A Simplified Explanation of the Transformer Block}
\item \href{https://arxiv.org/pdf/1706.03762}{Attention is All You Need}
\end{itemize}

Qwen3 uses the following Transformer block structure:

\begin{zhpar}
Qwen3 使用如下的 Transformer 块结构：
\end{zhpar}
\begin{codeblock}
  input
/ |
| input_layernorm (RMSNorm)
| |
| Qwen3MultiHeadAttention
\ |
  Add (residual)
/ |
| post_attention_layernorm (RMSNorm)
| |
| MLP
\ |
  Add (residual)
  |
output
\end{codeblock}

Run the tests for this task with:

\begin{zhpar}
用以下命令运行本任务的测试：
\end{zhpar}
\begin{codeblock}
pdm run test --week 1 --day 5 -- -k task_1
\end{codeblock}

\subsection*{Task 2: Implement \texttt{Embedding}}

\subsection*{任务 2：实现 \texttt{Embedding}}

\begin{codeblock}
src/tiny_llm/embedding.py
\end{codeblock}

\textbf{ 延伸阅读 / Readings}

\begin{itemize}[leftmargin=1.6em,itemsep=0.15em,parsep=0pt]
\item \href{https://huggingface.co/spaces/hesamation/primer-llm-embedding}{LLM Embeddings Explained: A Visual and Intuitive Guide}
\end{itemize}

The embedding layer maps token IDs (integers) to vectors of length \texttt{embedding\_dim}. In this task, you will implement that lookup operation.

\begin{zhpar}
嵌入层把 token ID（整数）映射为长度为 \texttt{embedding\_dim} 的向量。本任务中，你将实现这一查表操作。
\end{zhpar}
\begin{codeblock}
Embedding::__call__
weight: vocab_size x embedding_dim
Input: N.. (tokens)
Output: N.. x embedding_dim (vectors)
\end{codeblock}

This can be implemented with array indexing.

\begin{zhpar}
这可以用数组索引来实现。
\end{zhpar}
When input and output embeddings are tied, Qwen3 also uses the embedding weight as a linear projection from hidden vectors back to vocabulary logits.

\begin{zhpar}
当输入与输出嵌入\textbf{权重绑定（tied）}时，Qwen3 会直接复用嵌入权重，作为把隐藏向量映射回词表 logits 的线性投影。
\end{zhpar}
\begin{codeblock}
Embedding::as_linear
weight: vocab_size x embedding_dim
Input: N.. x embedding_dim
Output: N.. x vocab_size
\end{codeblock}

Run the tests for this task with:

\begin{zhpar}
用以下命令运行本任务的测试：
\end{zhpar}
\begin{codeblock}
# This task's tests use the 0.6B model and tokenizer.
hf download Qwen/Qwen3-0.6B-MLX-4bit
pdm run test --week 1 --day 5 -- -k task_2
\end{codeblock}

\subsection*{Task 3: Implement \texttt{Qwen3ModelWeek1}}

\subsection*{任务 3：实现 \texttt{Qwen3ModelWeek1}}

Now that we have built all the Qwen3 components, we can implement \texttt{Qwen3ModelWeek1}.

\begin{zhpar}
既然已经实现了 Qwen3 的全部组件，现在可以实现 \texttt{Qwen3ModelWeek1} 了。
\end{zhpar}
\begin{codeblock}
src/tiny_llm/qwen3_week1.py
\end{codeblock}

You will not implement the process of reading model parameters from tensor files. Instead, load the model with \texttt{mlx\_lm}, then transfer its parameters into our implementation. The \texttt{Qwen3ModelWeek1} constructor therefore accepts an MLX model.

\begin{zhpar}
你不需要实现“从张量文件读取模型参数”的过程。取而代之，用 \texttt{mlx\_lm} 加载模型，再把它的参数搬运进我们的实现中。因此 \texttt{Qwen3ModelWeek1} 的构造函数接收一个 MLX 模型对象。
\end{zhpar}
The Qwen3 model has the following layers:

\begin{zhpar}
Qwen3 模型包含如下各层：
\end{zhpar}
\begin{codeblock}
input
| (tokens: N..)
Embedding
| (N.. x hidden_size); note that hidden_size == embedding_dim
Qwen3TransformerBlock
| (N.. x hidden_size)
Qwen3TransformerBlock
| (N.. x hidden_size)
...
|
RMSNorm 
| (N.. x hidden_size)
Embedding.as_linear OR linear (lm_head)
| (N.. x vocab_size)
output
\end{codeblock}

Read the number of layers, hidden size, head dimension, and other configuration values from \texttt{mlx\_model.args}, whose type is defined by \href{https://github.com/ml-explore/mlx-lm/blob/main/mlx_lm/models/qwen3.py}{\texttt{ModelArgs}}. The loaded weights are available through \texttt{mlx\_model.model}; use the Qwen3 implementation and model metadata to identify the corresponding layer names.

\begin{zhpar}
层数、hidden size、head dimension 以及其他配置值，都从 \texttt{mlx\_model.args} 中读取，其类型由 \href{https://github.com/ml-explore/mlx-lm/blob/main/mlx_lm/models/qwen3.py}{\texttt{ModelArgs}} 定义。加载好的权重可以通过 \texttt{mlx\_model.model} 访问；请结合 Qwen3 的实现与模型元数据来确定各层对应的名称。
\end{zhpar}
By this point, you have implemented \texttt{RMSNorm}. Replace the temporary Day 3 calls to \texttt{mx.fast.rms\_norm} with \texttt{RMSNorm(head\_dim, q\_norm, eps=...)} and \texttt{RMSNorm(head\_dim, k\_norm, eps=...)}. They implement the same formula; the built-in calls existed only to keep the GQA chapter focused on attention.

\begin{zhpar}
到这里，你已经实现了 \texttt{RMSNorm}。请把第 3 天临时调用的 \texttt{mx.fast.rms\_norm} 替换为 \texttt{RMSNorm(head\_dim, q\_norm, eps=...)} 和 \texttt{RMSNorm(head\_dim, k\_norm, eps=...)}。二者实现的是同一个公式；当初使用内置调用，只是为了让 GQA 一章专注在注意力上。
\end{zhpar}
Different Qwen3 model variants map hidden vectors back to vocabulary logits in different ways. Some tie the input and output embeddings and use \texttt{Embedding.as\_linear}; others have a separate \texttt{lm\_head} linear layer. Select the strategy with \texttt{mlx\_model.args.tie\_word\_embeddings}: if it is \texttt{True}, use \texttt{Embedding.as\_linear}; otherwise, load and use \texttt{lm\_head}.

\begin{zhpar}
不同的 Qwen3 模型变体，把隐藏向量映射回词表 logits 的方式并不相同：有的绑定输入/输出嵌入并使用 \texttt{Embedding.as\_linear}，有的则带有独立的 \texttt{lm\_head} 线性层。请依据 \texttt{mlx\_model.args.tie\_word\_embeddings} 来选择策略：若为 \texttt{True}，使用 \texttt{Embedding.as\_linear}；否则加载并使用 \texttt{lm\_head}。
\end{zhpar}
The model takes a sequence of token IDs and returns unnormalized logits for every sequence position. On Day 6, we will use the final position's logits to select the next token and generate a response.

\begin{zhpar}
该模型接收一串 token ID，并返回每个序列位置上的未归一化 logits（unnormalized logits）。第 6 天，我们将使用最后一个位置的 logits 来选出下一个 token，从而生成回答。
\end{zhpar}
The MLX models used in this course have quantized weights. Dequantize each linear or embedding layer before loading it into tiny-llm by using the provided \texttt{quantize.dequantize\_linear} function, then store the readable Week 1 weight as BF16. Model activations and layer outputs should remain BF16. A readable attention or normalization expression may compute in FP32 for stability, but it must cast its model-facing result back to BF16.

\begin{zhpar}
本课程使用的 MLX 模型权重是量化过的。在把每个线性层或嵌入层加载进 tiny-llm 之前，请先使用已提供的 \texttt{quantize.dequantize\_linear} 函数对其做反量化，然后把这份“可读的 Week 1 权重”存成 BF16。模型的激活值与各层输出应保持 BF16。为了数值稳定，可读性优先的注意力或归一化表达式可以在 FP32 下计算，但必须把最终对外输出的结果转回 BF16。
\end{zhpar}
Pass \texttt{mask="causal"} to every Transformer block. For a one-token sequence the mask has no effect; for longer sequences, it prevents each position from attending to future tokens.

\begin{zhpar}
向每个 Transformer 块传入 \texttt{mask="causal"}。对单 token 序列而言该掩码不起作用；对更长的序列，它能阻止每个位置去 attend 未来的 token。
\end{zhpar}
Run the tests for this task with:

\begin{zhpar}
用以下命令运行本任务的测试：
\end{zhpar}
\begin{codeblock}
# Download each model you want to test. Missing models are skipped.
hf download Qwen/Qwen3-0.6B-MLX-4bit
hf download Qwen/Qwen3-1.7B-MLX-4bit
hf download Qwen/Qwen3-4B-MLX-4bit
pdm run test --week 1 --day 5 -- -k task_3
\end{codeblock}

At the end of the day, you should be able to pass all tests of this day:

\begin{zhpar}
在今天结束时，你应该能通过当天的全部测试：
\end{zhpar}
\begin{codeblock}
pdm run test --week 1 --day 5
\end{codeblock}

\vspace{1em}\begin{center}\small\sffamily tiny-llm-book \textcopyright{} 2025 by Alex Chi Z, CC BY-NC-SA 4.0. 英文原文 \url{https://github.com/skyzh/tiny-llm}；中文译文为对照学习用途。\end{center}
\end{document}
